\chapter{CÁC PROMPT SỬ DỤNG VỚI MÔ HÌNH NGÔN NGỮ}
\label{appendix:C}

Phụ lục này tập hợp các prompt dùng với mô hình ngôn ngữ trong hai vai trò: (1) vận hành hệ thống hỏi đáp (định tuyến ý định và sinh câu trả lời), và (2) xây dựng dữ liệu huấn luyện (xác định positive, sinh query-positive bổ sung). Các trường trong dấu ngoặc nhọn như \plaincode{\{question\}} là chỗ điền nội dung động khi chạy.

\section{Prompt định tuyến ý định}
\label{appendix:C.1}

Prompt phân loại câu hỏi thành một trong hai lớp \plaincode{tra_cuu} hoặc \plaincode{tinh_toan}.

\begin{promptbox}
Phân loại intent của câu hỏi sau. Trả lời chỉ một từ:
- "tra_cuu" nếu câu hỏi yêu cầu tra cứu, tìm kiếm thông tin, định nghĩa, giải thích
- "tinh_toan" nếu câu hỏi yêu cầu tính toán, suy luận số liệu, so sánh theo phép tính

Câu hỏi: {question}

Intent:
\end{promptbox}

\section{Prompt sinh câu trả lời theo ý định}
\label{appendix:C.2}

Với câu \plaincode{tra_cuu}, hệ thống dùng một lượt gọi để chọn đáp án trực tiếp:

\begin{promptbox}
Dựa vào thông tin sau đây, trả lời câu hỏi trắc nghiệm.

CONTEXT:
{context}

QUESTION:
{question}

OPTIONS:
{options_text}
Chỉ trả lời đúng 1 dòng, chỉ ghi chữ cái đáp án (A/B/C/D).
Nếu có nhiều đáp án đúng thì ghi tất cả, ví dụ: AB hoặc ACD.
Không giải thích.
ANSWER:
\end{promptbox}

Với câu \plaincode{tinh_toan}, hệ thống dùng hai lượt gọi. Lượt một sinh phân tích từng bước:

\begin{promptbox}
Dựa vào context sau, phân tích câu hỏi từng bước.

CONTEXT:
{context}

QUESTION:
{question}

OPTIONS:
{options_text}
Phân tích từng bước, giữ nguyên số liệu quan trọng.
Chưa cần kết luận đáp án cuối cùng:
\end{promptbox}

Lượt hai chọn đáp án cuối từ phần phân tích:

\begin{promptbox}
Dựa vào phân tích sau, chọn đáp án đúng cho câu hỏi trắc nghiệm.

PHÂN TÍCH:
{reasoning}

QUESTION:
{question}

OPTIONS:
{options_text}
Chỉ trả lời đúng 1 dòng, chỉ ghi chữ cái đáp án (A/B/C/D).
Nếu có nhiều đáp án đúng thì ghi tất cả, ví dụ: AB hoặc ACD.
Không giải thích.
ANSWER:
\end{promptbox}

\section{Prompt giám khảo xác định positive}
\label{appendix:C.3}

Prompt yêu cầu mô hình phân loại từng đoạn ứng viên thành gold, partial hoặc irrelevant, trả về kết quả dạng JSON.

\begin{promptbox}
Judge relevance cho RAG tiếng Việt.

Phân loại từng chunk:
- gold:       đủ thông tin trả lời trực tiếp
- partial:    liên quan nhưng không đủ
- irrelevant: không liên quan

JSON only:
{"gold_indices": [...], "partial_indices": [...],
 "irrelevant_indices": [...], "confidence": "high|medium|low"}
\end{promptbox}

Phần prompt người dùng kèm câu hỏi và danh sách đoạn:

\begin{promptbox}
QUERY: {question}
INTENT: {intent}

CÁC ĐOẠN VĂN:
[CHUNK_0] (bge_score={score})
{chunk_text}
...

Phân loại từng chunk.
\end{promptbox}

\section{Prompt sinh query-positive bổ sung}
\label{appendix:C.4}

Prompt này được dùng để tạo thêm câu hỏi cho các đoạn chưa xuất hiện trong tập huấn luyện hoặc đánh giá. System prompt đặt vai trò và yêu cầu trả về JSON hợp lệ:

\begin{promptbox}
Bạn tạo dữ liệu query-positive cho hệ thống truy hồi tài liệu.

Nhiệm vụ: đọc một đoạn văn pháp lý/kỹ thuật và viết các câu hỏi mà người dùng thật có thể gõ để tìm đoạn đó.
Luôn trả về JSON hợp lệ, không giải thích thêm.
\end{promptbox}

User prompt nhận nội dung đoạn và số câu hỏi cần sinh:

\begin{promptbox}
Bạn là người dùng đang tra cứu tài liệu pháp lý/kỹ thuật. Đọc đoạn văn sau và đặt câu hỏi mà đoạn này trả lời được.

ĐOẠN VĂN:
{chunk_text}

YÊU CẦU:
- Đặt {n_questions} câu hỏi KHÁC NHAU mà đoạn văn trên trả lời được.
- Viết như người dùng thật gõ tìm kiếm: ngắn gọn, tự nhiên. KHÔNG "xin vui lòng", KHÔNG trang trọng.
- QUAN TRỌNG: KHÔNG sao chép nguyên cụm từ dài trong đoạn văn. Diễn đạt lại bằng từ ngữ của bạn, dùng từ đồng nghĩa. Câu hỏi và đoạn văn KHÔNG được trùng lặp nhiều từ khóa.
- Câu hỏi phải trả lời được CHỈ bằng đoạn văn này, không cần thông tin ngoài.
- Nếu đoạn văn có số liệu/công thức/điều kiện tính toán, đặt ít nhất 1 câu hỏi dạng tính toán/áp dụng.
- Nếu đoạn văn quá ngắn/không có nội dung đáng hỏi (mục lục, tiêu đề, bảng trống), trả về danh sách rỗng.

Trả về JSON:
{"questions": ["câu hỏi 1", "câu hỏi 2"], "chunk_quality": "good" | "low"}
\end{promptbox}
